\documentclass[11pt]{article}
\usepackage[margin=2.1cm]{geometry}
\usepackage{booktabs}
\usepackage{amsmath}
\usepackage{amssymb}
\usepackage{enumitem}
\usepackage{array}
\usepackage{listings}
\usepackage{xcolor}
\usepackage[T1]{fontenc}
\usepackage{lmodern}

\setlength{\parindent}{0pt}
\setlength{\parskip}{3.6pt}
\newcommand{\code}[1]{\texttt{#1}}
\newcommand{\hr}{\vspace{2pt}\hrule\vspace{6pt}}
\renewcommand{\arraystretch}{1.12}

\definecolor{kw}{RGB}{0,90,160}
\definecolor{cmt}{RGB}{110,110,110}
\lstdefinestyle{vhdl}{
  language=VHDL, basicstyle=\ttfamily\footnotesize,
  keywordstyle=\color{kw}\bfseries, commentstyle=\color{cmt}\itshape,
  showstringspaces=false, columns=fullflexible, breaklines=true,
  frame=leftline, framesep=6pt, xleftmargin=8pt, aboveskip=6pt, belowskip=6pt
}
\lstset{style=vhdl}

\begin{document}

\begin{center}
{\Large\bfseries EE\,362 --- Pre-Lab 6: Solution}\\[2pt]
{\large Register File}\\[4pt]
Fall 2026 \quad$\cdot$\quad 5 points \quad$\cdot$\quad TA marking guide
\end{center}
\hr

{\small Marking: Q1 = 2 points, Q2 = 1 point, Q3 = 2 points.}

\vspace{4pt}\hr

\section*{Question 1 --- The write behaviour (2 points)}

\subsection*{(a) The table}

\begin{center}
\begin{tabular}{c|cc|l}
\toprule
\# & \code{Strb} & \code{Dsel} & effect at the rising edge \\
\midrule
0 & 0 & --- & \textbf{nothing is written}; every register holds \\
1 & 1 & \code{"00000"} & \textbf{nothing is written}: register 0 cannot be changed \\
2 & 1 & $k \neq 0$ & register $k$ takes \code{Din}; all others hold \\
\bottomrule
\end{tabular}
\end{center}

\textbf{Row 1 is the one being marked.} ``Writes zero'' or ``writes \code{Din} to
register 0'' both earn no credit --- register 0 is not \emph{cleared} by a write,
it is \emph{unaffected} by one.

\subsection*{(b) The write-enable equation}

\[
\mathit{We}_k = \code{Strb} \cdot m_k(\code{Dsel}) \cdot \boxed{[\,k \neq 0\,]}
\]

The boxed term is the one to circle: it is present only for $k = 0$, where it is
constant 0, so $\mathit{We}_0$ can never be asserted.

\begin{quote}
\textbf{What ``hardwired to zero'' means in hardware:} register 0 is not a
register that happens to contain zero. Its write enable is permanently false and
its read path is forced to zero, so \textbf{there is no storage for register 0 at
all} --- a good synthesiser optimises those 32 flip-flops away entirely.
\end{quote}

Full credit needs the idea that it cannot be written \emph{and} that reading it
always yields zero. A student who protects only the write side has a design that
reads back whatever the register file was initialised with.

\textbf{Note the shape of the equation}: $m_k(\code{Dsel})$ is a minterm of a
5-bit address --- exactly the decoder from Lab 0, Example 3, at 5-to-32.

\subsection*{(c) Three components already built}

\begin{center}
\begin{tabular}{lll}
\toprule
component & from & role here \\
\midrule
\textbf{decoder} (5-to-32) & Lab 0, Example 3 & turns \code{Dsel} into 32 write enables \\
\textbf{register} (\code{RegN}) & Lab 2 & 32 copies, each 32 bits wide \\
\textbf{multiplexer} (32-to-1) & Lab 3 & each read port selects one register \\
\bottomrule
\end{tabular}
\end{center}

\begin{quote}
\textbf{The point of the question:} nothing in this lab is new. A register file is
a decoder, a pile of registers, and two multiplexers. Students who see that build
it in an hour; students who see it as a new kind of thing do not.
\end{quote}

\vspace{4pt}\hr

\section*{Question 2 --- Interface and cost (1 point)}

\subsection*{(a) The entity}

\begin{lstlisting}
entity RegFile is
    port (
        Asel : in  std_logic_vector(4 downto 0);   -- read port A address
        Bsel : in  std_logic_vector(4 downto 0);   -- read port B address
        Dsel : in  std_logic_vector(4 downto 0);   -- write address
        Din  : in  std_logic_vector(31 downto 0);
        Strb : in  std_logic;                      -- write strobe
        clk  : in  std_logic;
        Aout : out std_logic_vector(31 downto 0);
        Bout : out std_logic_vector(31 downto 0)
    );
end entity RegFile;
\end{lstlisting}

\textbf{Five bits per address}, because $2^5 = 32$. A student who wrote 32-bit
addresses has confused the \emph{number} of registers with their \emph{width} ---
the same narrow-control / wide-data distinction as every lab since Lab 1.

\subsection*{(b) Why a third read port is cheap and a second write port is not}

\begin{quote}
\textbf{Model answer:} A read port is just another multiplexer looking at the
same stored values. Adding one costs one more 32-to-1 mux and changes nothing
else, because reading does not modify anything and any number of readers can look
at the same register at once.

A second write port means two sources can target the \textbf{same register on the
same edge}. The hardware must then decide which one wins, so it needs arbitration
or priority logic, and every register's input becomes a multiplexer between the
two write buses. The cost is in the \emph{conflict}, not in the wires.
\end{quote}

Full credit for ``readers cannot conflict, writers can''. A good extra sentence:
this is why the MIPS instruction format has two source fields and only one
destination field.

\vspace{4pt}\hr

\section*{Question 3 --- Test plan (2 points)}

\subsection*{Model vector list}

\begin{center}
\small
\begin{tabular}{c>{\raggedright\arraybackslash}p{5.4cm}>{\raggedright\arraybackslash}p{7.2cm}}
\toprule
\# & what you do & what it proves \\
\midrule
1 & write a pattern to register 5, then read it on port A & the basic write-then-read path works \\
2 & read the same register on port B & \textbf{both} read ports work, and independently \\
3 & read two \emph{different} registers on A and B at once & the ports do not share a selector \\
4 & write with \code{Strb} = 0 & no write happens without the strobe \\
5 & \textbf{write to register 0, then read it} & \textbf{register 0 is protected} \\
6 & write register $k$, then read a \emph{different} register $j$ & the decoder does not enable neighbours \\
7 & read a register while writing it on the same edge & read-during-write behaviour (see below) \\
\bottomrule
\end{tabular}
\end{center}

Six vectors earn full credit. \textbf{Vectors 5 and 6 are the discriminating
ones.}

Vector 6 deserves a comment: a decoder bug usually enables \emph{two} registers,
not none, so the way you see it is that some other register changed. Writing
distinct patterns to several registers and reading them all back is the cheap way
to catch it.

\subsection*{(a) Which vector catches an unprotected register 0?}

\begin{quote}
\textbf{Model answer:} Vector 5 --- write a non-zero pattern to register 0 with
\code{Strb} high, then read register 0. A correct file returns zero; an
unprotected one returns the pattern.
\end{quote}

\textbf{The vector must use a non-zero pattern.} Writing zeros to register 0 and
reading zero back proves nothing at all.

\subsection*{(b) Read-during-write: two correct behaviours}

\begin{center}
\begin{tabular}{ll}
\toprule
behaviour & what the read port returns \\
\midrule
\textbf{read-old} & the value stored \emph{before} the edge \\
\textbf{read-new} (write-through, bypass) & the value being written \emph{on} that edge \\
\bottomrule
\end{tabular}
\end{center}

\begin{quote}
\textbf{Model answer:} Both are legitimate designs; they are different
specifications, not a right and a wrong one. To tell them apart, put a known old
value in register $k$, then on one edge write a \emph{different} value to $k$
while reading $k$, and sample the read port \textbf{just before} the edge and
again just after. Read-old shows the previous value up to and including that
edge's read; read-new shows the new value immediately.
\end{quote}

\textbf{This course uses read-old}, because the read ports are combinational
multiplexers looking at stored outputs, and nothing is stored until the edge. No
bypass path is built.

\textbf{The consequence, worth saying now:} it means an instruction cannot read a
value written by the instruction immediately before it in the same cycle. In this
single-cycle CPU that never happens, because each instruction has a whole cycle.
It is exactly the hazard that appears in the Final Lab's pipeline track.

\vspace{6pt}
\hr
{\small \textbf{At the bench:} the failure to look for is a student who protected
register 0 on the write side only. Their vector 5 passes if the array was
initialised to zeros, and fails as soon as anything ever targets register 0.}

\end{document}
